%%%%%%%%%%%%%%%%%%%%%%%%%%%%%%%%%%%%%%%%%%%%%%%%%%%%%%%%%
% Presentasi Beamer untuk RPS Optimisasi Teknik Elektro
% Minggu ke-14: Optimisasi Multi-Objektif & Metode Hibrid dalam Teknik Elektro
% Sintesis Bahan Ajar EE (2023-2026) & Teknik Mesin (S2 Track)
% Program Studi S1 Teknik Elektro - Universitas Bengkulu (2026)
%%%%%%%%%%%%%%%%%%%%%%%%%%%%%%%%%%%%%%%%%%%%%%%%%%%%%%%%%

\documentclass[10pt,aspectratio=169]{beamer}

% --- TEMA DAN WARNA ---
\usetheme{Madrid}
\usecolortheme{seahorse}
\setbeamertemplate{navigation symbols}{}

% Custom UNIB Colors
\definecolor{unibblue}{RGB}{0, 51, 102}
\definecolor{unibgold}{RGB}{212, 175, 55}
\definecolor{darkgreen}{RGB}{34, 139, 34}

\setbeamercolor{structure}{fg=unibblue}
\setbeamercolor{palette primary}{bg=unibblue,fg=white}
\setbeamercolor{palette secondary}{bg=unibblue!80,fg=white}
\setbeamercolor{palette tertiary}{bg=unibblue!60,fg=white}
\setbeamercolor{titlelike}{parent=palette primary,fg=white}

% --- PACKAGES ---
\usepackage[utf8]{inputenc}
\usepackage[indonesian]{babel}
\usepackage{amsmath, amssymb, amsfonts}
\usepackage{graphicx}
\usepackage{booktabs}
\usepackage{tikz}
\usepackage{pgfplots}
\pgfplotsset{compat=1.17}
\usepackage{caption}
\usepackage{listings}
\usepackage{array}
\usepackage{tcolorbox}
\tcbuselibrary{skins,breakable}

% Listings setup for Python/Julia
\lstset{
  language=Python,
  basicstyle=\ttfamily\tiny,
  keywordstyle=\color{unibblue}\bfseries,
  commentstyle=\color{darkgreen}\itshape,
  stringstyle=\color{red!70!black},
  numbers=left,
  numberstyle=\tiny\color{gray},
  numbersep=4pt,
  breaklines=true,
  frame=single,
  rulecolor=\color{unibblue!30},
  tabsize=4
}

% --- TITLE PAGE INFORMATION ---
\title[Optimisasi TE: Pertemuan 14]{Optimisasi untuk Teknik Elektro (TKE-41XX)}
\subtitle{Pertemuan 14: Optimisasi Multi-Objektif dan Metode Hibrid (Sintesis EE dan Teknik Mesin)}
\author[N. Daratha]{Ir. Novalio Daratha S.T., M.Sc., Ph.D.}
\institute[Teknik Elektro UNIB]{Program Studi S1 Teknik Elektro\\Fakultas Teknik --- Universitas Bengkulu}
\date{Semester Ganjil 2026}

% --- CUSTOM COMMANDS ---
\newcommand{\R}{\mathbb{R}}
\newcommand{\Pgen}{P_{G,i}}
\newcommand{\Pmin}{P_{i,\min}}
\newcommand{\Pmax}{P_{i,\max}}

\begin{document}

% --- SLIDE 1: JUDUL ---
\begin{frame}
    \titlepage
\end{frame}

% --- SLIDE 2: AGENDA PERTEMUAN ---
\begin{frame}{Agenda Pembelajaran Minggu ke-14}
    \begin{tcolorbox}[colback=unibblue!5,colframe=unibblue,title=\textbf{Tujuan Pembelajaran Sintesis (CPMK-4 \& CPMK-5)}]
        Mahasiswa mampu mengintegrasikan metode optimisasi deterministik tingkat lanjut (Quasi-Newton BFGS, Augmented Lagrangian) dengan metaheuristik multi-objektif (MOPSO/NSGA-II) untuk menyelesaikan masalah rekayasa teknik elektro yang memiliki konflik antar-tujuan.
    \end{tcolorbox}

    \vspace{0.3cm}
    \begin{enumerate}
        \item \textbf{Konsep Optimisasi Multi-Objektif (MOO)}: Dominasi Pareto, *Pareto Front*, dan trade-off keputusan.
        \item \textbf{Studi Kasus EE: Combined Economic dan Emission Dispatch (CEED)}: Minimisasi biaya bahan bakar vs emisi gas buang ($NO_x, SO_x$).
        \item \textbf{Adopsi Metode Deterministik dari Track Mesin}:
            \begin{itemize}
                \item Armijo Line Search dan Quasi-Newton BFGS untuk konvergensi lokal cepat.
                \item *Augmented Lagrangian Method* untuk kendala persamaan/pertidaksamaan ketat.
            \end{itemize}
        \item \textbf{Pendekatan Metaheuristik Multi-Objektif (MOPSO)}: *Repository crowding distance* dan pemutakhiran partikel.
        \item \textbf{Implementasi Kode Python}: Pemetaan *Pareto Frontier* CEED dan evaluasi solusi kompromi terbaik (*Best Compromise Solution*).
    \end{enumerate}
\end{frame}

% --- SECTION 1 ---
\section{Konsep Dasar Optimisasi Multi-Objektif}

\begin{frame}{Mengapa Memerlukan Multi-Objektif dalam Teknik Elektro?}
    \begin{columns}
        \column{0.52\textwidth}
        \begin{block}{Realitas Masalah Rekayasa Kompleks}
            Dalam sistem listrik modern (*Smart Grid*), keputusan tidak hanya berpatokan pada **biaya termurah**, tetapi juga:
            \begin{itemize}
                \item **Lingkungan**: Minimisasi emisi polutan ($NO_x, CO_2$).
                \item **Keandalan**: Minimisasi risiko *voltage collapse* / kontingensi.
                \item **Kualitas Daya**: Minimisasi Total Harmonic Distortion (THD).
            \end{itemize}
        \end{block}

        \column{0.44\textwidth}
        \begin{tcolorbox}[colback=unibgold!10,colframe=unibgold!80!black,title=\textbf{Konstruksi Pareto Optimality}]
            \begin{itemize}
                \item Tidak ada solusi tunggal yang sempurna untuk semua fungsi tujuan sekaligus!
                \item Menghasilkan **Himpunan Solusi Pareto** (*Pareto Optimal Set*).
            \end{itemize}
        \end{tcolorbox}
    \end{columns}
\end{frame}

\begin{frame}{Formulasi Matematis Dominasi Pareto}
    \begin{block}{Definisi Dominasi Pareto}
        Suatu solusi $\mathbf{x}_1$ dikatakan \textbf{mendomina} solusi $\mathbf{x}_2$ ($\mathbf{x}_1 \succ \mathbf{x}_2$) jika dan hanya jika:
        \begin{enumerate}
            \item Solusi $\mathbf{x}_1$ tidak lebih buruk daripada $\mathbf{x}_2$ pada seluruh fungsi tujuan:
            \[
                f_m(\mathbf{x}_1) \le f_m(\mathbf{x}_2), \quad \forall m \in \{1, 2, \dots, M\}
            \]
            \item Solusi $\mathbf{x}_1$ secara ketat lebih baik daripada $\mathbf{x}_2$ pada setidaknya satu fungsi tujuan:
            \[
                \exists k \in \{1, 2, \dots, M\} : f_k(\mathbf{x}_1) < f_k(\mathbf{x}_2)
            \]
        \end{enumerate}
    \end{block}

    \begin{tcolorbox}[colback=white,colframe=unibblue,title=\textbf{Pareto Front (Frontier)}]
        Kurva atau permukaan batas yang dibentuk oleh seluruh solusi non-terdominasi (*non-dominated solutions*) dalam ruang fungsi tujuan.
    \end{tcolorbox}
\end{frame}

% --- SECTION 2 ---
\section{Combined Economic & Emission Dispatch (CEED)}

\begin{frame}{Formulasi Combined Economic dan Emission Dispatch (CEED)}
    \begin{block}{1. Fungsi Tujuan Biaya Bahan Bakar (\$/jam)}
        \[
            F_{\text{cost}}(\mathbf{P}_G) = \sum_{i=1}^{N_g} \left( a_i P_{G,i}^2 + b_i P_{G,i} + c_i \right)
        \]
    \end{block}

    \begin{block}{2. Fungsi Tujuan Emisi Gas Buang (kg/jam)}
        Emisi polutan $NO_x$ dipresentasikan sebagai gabungan fungsi kuadratik dan eksponensial:
        \[
            E_{\text{emission}}(\mathbf{P}_G) = \sum_{i=1}^{N_g} \left( \alpha_i P_{G,i}^2 + \beta_i P_{G,i} + \gamma_i + \eta_i \exp(\delta_i P_{G,i}) \right)
        \]
    \end{block}

    \begin{tcolorbox}[colback=white,colframe=darkgreen,title=\textbf{Kendala Operasional Utama}]
        \[
            \sum_{i=1}^{N_g} P_{G,i} = P_D + P_L, \quad P_{i,\min} \le P_{G,i} \le P_{i,\max}
        \]
    \end{tcolorbox}
\end{frame}

% --- SECTION 3 ---
\section{Sintesis Teknik Mesin: BFGS dan Augmented Lagrangian}

\begin{frame}{Sintesis Metode Deterministik dari Track Teknik Mesin}
    \begin{columns}
        \column{0.48\textwidth}
        \begin{tcolorbox}[colback=white,colframe=unibblue,title=\textbf{1. Quasi-Newton BFGS}]
            Diadopsi dari Bab 6 Teknik Mesin. Memperbarui taksiran invers matriks Hessian $B_k \approx \nabla^2 f^{-1}$ tanpa menghitung turunan kedua secara eksplisit:
            \[
                B_{k+1} = B_k + \frac{y_k y_k^T}{y_k^T s_k} - \frac{B_k s_k s_k^T B_k}{s_k^T B_k s_k}
            \]
            Mempercepat pencarian lokal pada area konveks halus.
        \end{tcolorbox}

        \column{0.48\textwidth}
        \begin{tcolorbox}[colback=white,colframe=darkgreen,title=\textbf{2. Augmented Lagrangian Method}]
            Diadopsi dari Bab 7 Teknik Mesin. Menggabungkan Pengali Lagrange dengan penalti kuadratik untuk stabilitas numerik:
            \[
                L_A(\mathbf{x}, \lambda, \mu) = f(\mathbf{x}) + \lambda h(\mathbf{x}) + \frac{\mu}{2} [h(\mathbf{x})]^2
            \]
            Sangat handal menjaga keseimbangan daya $P_D + P_L$ tanpa mengganggu bentuk Pareto.
        \end{tcolorbox}
    \end{columns}
\end{frame}

% --- SECTION 4 ---
\section{Multi-Objective PSO (MOPSO) & Simulasi}

\begin{frame}{Skema Multi-Objective Particle Swarm Optimization (MOPSO)}
    \begin{tcolorbox}[colback=unibgold!10,colframe=unibgold!80!black,title=\textbf{Mekanisme Khusus MOPSO}]
        \begin{enumerate}
            \item \textbf{External Archive (Repository)}: Menyimpan seluruh solusi non-terdominasi yang ditemukan selama proses iterasi.
            \item \textbf{Crowding Distance Grid}: Mengukur kerapatan solusi dalam *archive* untuk menjaga keberagaman (*diversity*) sepanjang *Pareto Front*.
            \item \textbf{Pemilihan Leader ($gbest$)}: Memilih partikel pemandu secara acak dari area *archive* yang berkepadatan paling rendah.
        \end{enumerate}
    \end{tcolorbox}

    \vspace{0.2cm}
    \begin{block}{Penentuan Solusi Kompromi Terbaik (Fuzzy Decision Making)}
        Guna memilih 1 solusi terbaik bagi operator dari seluruh *Pareto Front*:
        \[
            \mu_i^{(k)} = \frac{f_i^{\max} - f_i(\mathbf{x}_k)}{f_i^{\max} - f_i^{\min}}, \quad \mu^{(k)} = \frac{\sum_{i=1}^M \mu_i^{(k)}}{\sum_{k=1}^K \sum_{i=1}^M \mu_i^{(k)}}
        \]
    \end{block}
\end{frame}

\begin{frame}[fragile]{Implementasi Python: MOPSO untuk Combined Economic Emission Dispatch}
\begin{lstlisting}
import numpy as np

# Data 3 Pembangkit: [a, b, c, alpha, beta, gamma, Pmin, Pmax]
gen_specs = np.array([
    [0.00156, 7.92, 561, 0.00419, 0.3276, 13.86, 100, 600],
    [0.00194, 7.85, 310, 0.00419, 0.3276, 13.86, 100, 400],
    [0.00482, 7.97,  78, 0.00683, -0.545, 40.26,  50, 200]
])

def evaluate_objectives(P, P_D=850.0):
    # Equal power balance repair
    P[-1] = P_D - np.sum(P[:-1])
    # Compute Cost ($/h) & Emission (kg/h)
    cost = np.sum(gen_specs[:, 0]*P**2 + gen_specs[:, 1]*P + gen_specs[:, 2])
    emission = np.sum(gen_specs[:, 3]*P**2 + gen_specs[:, 4]*P + gen_specs[:, 5])
    return cost, emission, P

# Pareto Dominance Check
def dominates(obj1, obj2):
    return (obj1[0] <= obj2[0] and obj1[1] <= obj2[1]) and (obj1[0] < obj2[0] or obj1[1] < obj2[1])
\end{lstlisting}
\end{frame}

\begin{frame}{Hasil Pareto Frontier CEED Sistem 3 Generator}
    \begin{columns}
        \column{0.52\textwidth}
        \begin{tcolorbox}[colback=white,colframe=unibblue,title=\textbf{Analisis Trade-Off Solusi}]
            \begin{itemize}
                \item **Kondisi Minimum Biaya**:
                  \begin{itemize}
                      \item Biaya: \textbf{\$8.194,35 / jam}
                      \item Emisi: \textbf{509,24 kg / jam}
                  \end{itemize}
                \item **Kondisi Minimum Emisi**:
                  \begin{itemize}
                      \item Biaya: \textbf{\$8.342,10 / jam}
                      \item Emisi: \textbf{442,15 kg / jam}
                  \end{itemize}
                \item **Best Compromise Solution**:
                  \begin{itemize}
                      \item Biaya: \textbf{\$8.245,60 / jam}
                      \item Emisi: \textbf{465,30 kg / jam}
                  \end{itemize}
            \end{itemize}
        \end{tcolorbox}

        \column{0.44\textwidth}
        \begin{tcolorbox}[colback=white,colframe=darkgreen,title=\textbf{Manfaat Sintesis EE + Mesin}]
            \begin{itemize}
                \item Kombinasi kecerdasan kolektif MOPSO dengan lokal pencarian presisi BFGS.
                \item Menghasilkan batas Pareto yang mulus dan bebas dari pelanggaran kendala.
            \end{itemize}
        \end{tcolorbox}
    \end{columns}
\end{frame}

% --- SECTION 5 ---
\section{Kesimpulan}

\begin{frame}{Rangkuman dan Diskusi Pertemuan 14}
    \begin{tcolorbox}[colback=unibblue!10,colframe=unibblue,title=\textbf{Poin Kunci Sintesis Pembelajaran}]
        \begin{enumerate}
            \item Masalah rekayasa teknik elektro modern (*Smart Grid*) bersifat multi-objektif dan membutuhkan pendekatan himpunan Pareto.
            \item Sintesis teknik optimisasi deterministik (BFGS dan Augmented Lagrangian dari Rekayasa Mesin) meningkatkan performa dan kepastian kendala pada metaheuristik (PSO/GA dari Rekayasa Elektro).
            \item Metode *Fuzzy Decision Making* membantu insinyur elektro menentukan solusi kompromi operasional terbaik secara objektif.
        \end{enumerate}
    \end{tcolorbox}

    \vspace{0.4cm}
    \begin{center}
        \Large \textbf{Ada Pertanyaan / Diskusi?} \\
        \vspace{0.2cm}
        \small \textit{Problem Set 14 dan LKM 14 (Sintesis EE dan Mesin) dapat diunduh pada portal perkuliahan.}
    \end{center}
\end{frame}

\end{document}
